\exercisesection{代表域}

\begin{enumerate}
\def\labelenumi{\arabic{enumi})}
\tightlist
\item
  设 \(p\) 是素数，\(G\) 是以 \(X\)、\(Y\) 两个生成元生成的自由群，\(F_p[G]\) 是 \(G\) 在 \(F_p\) 上的群代数。置 \(Z = XY - YX\)，并用 \(\mathfrak{a}\) 表示 \(F_p[G]\) 的由元素 \(ZX - XZ\)、\(ZY - YZ\) 与 \(Z^2\) 生成的双边理想。置 \(A = F_p[G]/\mathfrak{a}\)，并用 \(\mathfrak{z}\) 表示 \(A\) 的由 \(Z\) 在 \(A\) 中的像生成的双边理想（我们把 \(X\)、\(Y\)、\(Z\) 在 \(A\) 中的像仍记作 \(X\)、\(Y\)、\(Z\)）。a) 设 \(a\) 是 \(A\) 的一个元素。证明存在两个由 \(F_p\) 的元素组成的有限支撑族，即 \((p_{\alpha,\beta})_{(\alpha,\beta)\in\mathbf{Z}^2}\) 与 \((q_{\alpha,\beta})_{(\alpha,\beta)\in\mathbf{Z}^2}\)，它们由条件
\end{enumerate}

\[
a = \sum_ {(\alpha , \beta) \in \mathbf {Z} ^ {2}} (p _ {\alpha , \beta} + q _ {\alpha , \beta} \mathbf {Z}) \mathrm{X} ^ {\alpha} \mathrm{Y} ^ {\beta}.
\]

唯一确定。由此推出商代数 \(A/\mathfrak{z}\) 是交换的且是整的，并且理想 \(\mathfrak{z}\) 包含在 A 的中心中。

\begin{enumerate}
\def\labelenumi{\alph{enumi})}
\setcounter{enumi}{1}
\tightlist
\item
  对 A 的元素的每个偶 \((a, b)\)，存在 A 的元素 c，使得
\end{enumerate}

\[
  a ^ {k} b - b a ^ {k} = k a ^ {k - 1} c Z \quad \text { for every integer } k \geqslant 0.
\]

由此推出，A 的元素的 p 次幂所成的集合 \(A^{p}\) 包含在 A 的中心中。

\begin{enumerate}
\def\labelenumi{\alph{enumi})}
\setcounter{enumi}{2}
\tightlist
\item
  乘性部分 \(S = A - \mathfrak{z}\)（\(A\) 的）允许在右边与在左边做分式（II, § 2, 习题 22）。用 B 表示这样构造的 \(A\) 的分式环。
\end{enumerate}

（例如，为验证对 A 中的每个 \(a\) 与 S 中的每个 \(s\)，存在 \(b \in A\) 与 \(t \in S\) 使得 \(at = sb\)，可取 \(b = s^{p-1}a\)，\(t = s^p\)，并利用 b)。）

\begin{enumerate}
\def\labelenumi{\alph{enumi})}
\setcounter{enumi}{3}
\item
  证明从 A 到 B 的典范映射是单射。（注意这个映射的核包含在 \(\mathfrak{z}\) 中，并按 a) 的方式推理。）
\item
  B 的中心包含 Z。若 m 是 B 的由 Z 生成的双边理想，则 \(m^{2} = 0\)。
\item
  理想 m 是 B 的不可逆元素的集合，且域 B/m 同构于域 \(\mathbf{F}_{p}(\mathbf{X}, \mathbf{Y})\)。
\item
  证明 B 中不存在其在 B/m 中的像为整个 B/m 的子域。
\end{enumerate}

在习题 2 至 5 中，假设（交换）拓扑环都是线性拓扑化的。若 A 与 B 是两个拓扑环，且 B 是 A-代数，则当定义代数结构的从 A 到 B 的同态为连续时，称 B 是拓扑 A-代数。

\begin{enumerate}
\def\labelenumi{\arabic{enumi})}
\setcounter{enumi}{1}
\tightlist
\item
  设 A 是拓扑环，B 是拓扑 A-代数。我们称 B 是形式光滑 A-代数\(^{1}\)，如果对每个离散拓扑 A-代数 C、C 的每个平方为零的理想 j 以及每个连续 A-同态 \(u: B \to C/j\)，存在通过过渡到商诱导 u 的连续 A-同态 \(v: B \to C\)。
\end{enumerate}

设 A 是环，B 是 A-代数。如果当 A 与 B 赋有离散拓扑时 B 是形式光滑 A-代数，则称 B 是光滑 A-代数。

\begin{enumerate}
\def\labelenumi{\alph{enumi})}
\item
  设 A 是拓扑环，B 是形式光滑 A-代数。设 C 是环，j 是 C 的理想。给 C 赋以 j-进拓扑，并假设 C 对这一拓扑是分离的且完备的。则对每个连续 A-同态 \(u:B\to C/\mathfrak{j}\)，存在连续 A-同态 \(v:B\to C\)，它通过过渡到商诱导 \(u\)。
\item
  设 A 是环。A 上的每个多项式代数都是光滑 A-代数。
\item
  设 A 是拓扑环。若 B 是形式光滑 A-代数，且 C 是形式光滑 B-代数，则 C 是形式光滑 A-代数。
\item
  设 A 是拓扑环，B 是形式光滑 A-代数，A' 是拓扑 A-代数。则拓扑 A'-代数 \(B \otimes_{A} A'\)（III, § 2, 习题 28）是形式光滑 A'-代数。
\item
  设 A 是拓扑环，B 是形式光滑 A-代数。设 S（相应地 T）是 A（相应地 B）的乘性子集，使得 S 在 B 中的像包含于 T。则 \(T^{-1}B\) 是形式光滑 \(S^{-1}A\)-代数（\(T^{-1}B\) 与 \(S^{-1}A\) 的拓扑见 III, § 2, 习题 27）。
\item
  设 n 是 \(\geqslant 1\) 的整数，\((\mathbf{B}_{i})_{1\leqslant i\leqslant n}\) 是拓扑 A-代数的一个族。要使 \(\prod_{i=1}^{n} B_{i}\) 是形式光滑 A-代数，必须且只需 \(B_{i}\) 对 \(1 \leqslant i \leqslant n\) 都是形式光滑 A-代数。
\item
  设 A 是拓扑环，B 是拓扑 A-代数，\(\hat{A}\) 与 \(\hat{B}\) 分别是 A 与 B 的分离完备化。则下列三个条件等价：
\end{enumerate}

\begin{enumerate}
\def\labelenumi{(\roman{enumi})}
\item
  B 是形式光滑 A-代数；
\item
  \(\hat{\mathbf{B}}\) 是形式光滑 A-代数；
\item
  \(\hat{\mathbf{B}}\) 是形式光滑 \(\hat{\mathbf{A}}\)-代数。
\end{enumerate}

\(h)\) 重新采用 \(e\)) 中的记号与假设。则 \(\mathbf{B}\{\mathrm{T}^{-1}\}\) 是形式光滑 \(\mathbf{A}\{\mathrm{S}^{-1}\}\)-代数。

\begin{enumerate}
\def\labelenumi{\arabic{enumi})}
\setcounter{enumi}{2}
\tightlist
\item
  设 \(k\) 是域，A 是交换 \(k\)-代数。对每个整数 \(i \geqslant 1\)，置 \(\mathbf{B}_i = \underbrace{\mathbf{A} \otimes_k \mathbf{A} \otimes_k \cdots \otimes_k \mathbf{A}}_{i \text{ 项}}\)，并给 \(\mathbf{B}_i\) 赋以让
\end{enumerate}

A 作用在第一个分量上所得到的 A-代数结构。置 \(C_{1} = B_{2}\)，\(C_{2} = B_{3}\)，\(C_{3} = B_{4} \oplus B_{3}\)，并用公式定义 A-线性映射 \(d_{2}:C_{2} \to C_{1}\) 与 \(d_{3}:C_{3} \to C_{2}\)

\[
d _ {3} ((1 \otimes x \otimes y \otimes z, 1 \otimes \alpha \otimes \beta)) = x \otimes y \otimes z - 1 \otimes x y \otimes z + 1 \otimes x \otimes y z - z \otimes x \otimes y + 1 \otimes \alpha \otimes \beta - 1 \otimes \beta \otimes \alpha,
\]

\[
d _ {2} (1 \otimes \alpha \otimes \beta) = \alpha \otimes \beta - 1 \otimes \alpha \beta + \beta \otimes \alpha ,
\]

\(^{1}\) 若 A 是（离散）域 k 上的局部的 Noether 的完备代数，则根据下面的习题 5，这一定义与第 VIII 卷第 98 页习题 30 中给出的定义一致。

对 A 的元素 \(x, y, z, \alpha, \beta\) 的任意选择，这样就得到一个 A-模的复形

\[
\mathrm{C} _ {k} (\mathrm{A}): \mathrm{C} _ {3} \xrightarrow {d _ {3}} \mathrm{C} _ {2} \xrightarrow {d _ {2}} \mathrm{C} _ {1}.
\]

若 N 是 A-模，则称 \(k\)-双线性映射 \(f: \mathbf{A} \times \mathbf{A} \to \mathbf{N}\) 为 2-上闭链，如果有恒等式 \(xf(y, z) - f(xy, z) + f(x, yz) - zf(x, y) = 0\)（对 A 中的 \(x, y, z\)）；称它是对称的，如果有 \(f(x, y) = f(y, x)\)（对每个偶 \((x, y) \in \mathbf{A} \times \mathbf{A}\)）。

\begin{enumerate}
\def\labelenumi{\alph{enumi})}
\tightlist
\item
  下列条件等价：
\end{enumerate}

\begin{enumerate}
\def\labelenumi{(\roman{enumi})}
\item
  复形 \(\operatorname{Hom}_{\mathbf{A}}(\mathbf{C}_k(\mathbf{A}), \mathbf{N})\) 对每个 A-模 \(\mathbf{N}\) 都是零调的；
\item
  对每个 A-模 N，每个对称 2-上闭链 \(f:A\times A\to\mathbf{N}\) 都是 1-上边缘，即形如 \(f(a,b)=ag(b)+bg(a)-g(ab)\)，其中 \(g\in\mathrm{Hom}_{k}(\mathbf{A},\mathbf{N})\)；
\item
  A 是 k 上的光滑代数。
\end{enumerate}

\begin{enumerate}
\def\labelenumi{\alph{enumi})}
\setcounter{enumi}{1}
\tightlist
\item
  若 A 是域，则前面的条件还等价于复形 \(C_{k}(A)\) 零调这一条件。
\end{enumerate}

\begin{enumerate}
\def\labelenumi{\arabic{enumi})}
\setcounter{enumi}{3}
\tightlist
\item
  \begin{enumerate}
  \def\labelenumii{\alph{enumii})}
  \tightlist
  \item
    设 \(k\) 是域，\(K\) 是 \(k\) 的一个扩张。若 \(K\) 在 \(k\) 上可分，则 \(K\) 是光滑 \(k\)-代数。（若 \(K\) 是 \(k\) 的有限生成扩张，利用 § 3, n° 2 命题 1。一般情形，把 \(K\) 写成 \(k\) 的有限生成扩张的并，并利用前一习题。）
  \end{enumerate}
\end{enumerate}

\begin{enumerate}
\def\labelenumi{\alph{enumi})}
\setcounter{enumi}{1}
\tightlist
\item
  设 A 是包含域 k 的分离且完备的局部环。利用 a)，证明 A 有一个表示域。若 A 的剩余域是 k 的可分扩张，则 A 有一个包含 k 的表示域。
\end{enumerate}

\begin{enumerate}
\def\labelenumi{\arabic{enumi})}
\setcounter{enumi}{4}
\tightlist
\item
  设 A 是包含域 \(k\) 的 Noether 局部环，m 是 A 的极大理想。假设 A 是形式光滑 \(k\)-代数（域 \(k\) 是离散的）。
\end{enumerate}

\begin{enumerate}
\def\labelenumi{\alph{enumi})}
\item
  设 K 是环 \(\mathbb{A}/\mathfrak{m}^2\) 的一个表示域（习题 4 b)），\(x_1, \ldots, x_d\) 是 m 中其像构成 K-向量空间 \(\mathfrak{m}/\mathfrak{m}^2\) 之基的元素。设 \(\mathrm{K}[\mathrm{X}_1, \ldots, \mathrm{X}_d]\) 是 \(d\) 个未定元 \(\mathrm{X}_1, \ldots, \mathrm{X}_d\) 的多项式环，\(\mathfrak{n}\) 是它由 \(\mathrm{X}_1, \ldots, \mathrm{X}_d\) 生成的理想，\(\varphi\) 是从 \(\mathrm{K}[\mathrm{X}_1, \ldots, \mathrm{X}_d]\) 到 \(\mathbb{A}/\mathfrak{m}^2\) 中的、把 \(\mathrm{X}_i\) 映到 \(x_i\)（\(1 \leqslant i \leqslant d\)） 且在 K 上诱导恒等映射的同态。则 \(\varphi\) 诱导同构 \(\overline{\varphi}\)（从 \(\mathrm{K}[\mathrm{X}_1, \ldots, \mathrm{X}_d]/\mathfrak{n}^2\) 到 \(\mathbb{A}/\mathfrak{m}^2\) 上的）。
\item
  对每个整数 \(n \geqslant 1\)，存在同态 \(\psi_n: \mathbb{A} \to \mathrm{K}[\mathrm{X}_1, \ldots, \mathrm{X}_d]/\mathfrak{n}^{n+1}\)，它通过过渡到商诱导 \(\overline{\varphi}^{-1}\)。这样的同态 \(\psi_n\) 是满射。
\item
  对每个整数 \(n \geqslant 1\)，A-模 \(\mathbf{A} / \mathfrak{m}^{n+1}\) 的长度至少为 \(\binom{d+n}{d}\)，且有 \(\dim(\mathbf{A}) \geqslant d\)。
\item
  对 \(k'\)（\(k\) 的每个有限扩张）以及每个极大理想 \(m'\)（半局部环 \(A \otimes_k k'\) 的），局部环 \((A \otimes_k k')_{m'}\) 都是正则的。
\end{enumerate}

\begin{enumerate}
\def\labelenumi{\arabic{enumi})}
\setcounter{enumi}{5}
\tightlist
\item
  设 \(k\) 是域，\(K\) 是 \(k\) 的一个扩张。假设 \(K\) 是光滑 \(k\)-代数。设 \(k'\) 是 \(k\) 的有限次数代数扩张。
\end{enumerate}

\begin{enumerate}
\def\labelenumi{\alph{enumi})}
\item
  环 K \(\otimes_{k} k'\) 是有限多个 Artin 局部环的乘积，这些局部环都是光滑 \(k'\)-代数（利用第 76 页习题 2）。
\item
  环 K \(\otimes_{k} k'\) 是既约的（利用前一习题）。
\item
  域 K 在 k 上是可分的。
\end{enumerate}

\begin{enumerate}
\def\labelenumi{\arabic{enumi})}
\setcounter{enumi}{6}
\tightlist
\item
  设 \(k\) 是域，\(P\) 是其素子域，\(K\) 是 \(k\) 的一个扩张。则下列条件等价：
\end{enumerate}

\begin{enumerate}
\def\labelenumi{\alph{enumi})}
\item
  k 到 K-模 M 中的每个导子都唯一地扩张为 K 到 M 中的导子。
\item
  我们有 \(\Omega_{\mathbf{P}}(\mathbf{K}) = \Omega_{\mathbf{P}}(k)\otimes_k\mathbf{K}\)。
\item
  域 K 在 k 上可分，且 \(\Omega_{k}(\mathbf{K}) = 0\)。
\end{enumerate}

若 \(k\) 的特征为 0，则这些条件还等价于条件：

\begin{enumerate}
\def\labelenumi{\alph{enumi})}
\setcounter{enumi}{3}
\tightlist
\item
  域 K 是 k 的代数扩张。
\end{enumerate}

若 \(k\) 的特征 \(p\ne0\)，则它们还等价于条件：

\begin{enumerate}
\def\labelenumi{\alph{enumi})}
\setcounter{enumi}{4}
\item
  \(\mathbf{K} = k\otimes_{k^p}\mathbf{K}^p.\)
\item
  k 的每个 p-基也是 K 的 p-基。
\end{enumerate}

\begin{enumerate}
\def\labelenumi{\arabic{enumi})}
\setcounter{enumi}{7}
\tightlist
\item
  设 \(k\) 是域，\(K\) 是 \(k\) 的一个扩张。我们称 \(K\) 在 \(k\) 上是形式平展的，如果对每个 \(k\)-代数 \(A\)、每个平方为零的理想 \(a\)（\(A\) 的）以及每个 \(k\)-同态 \(u:K\to A/a\)，存在唯一的 \(k\)-同态 \(v:K\to A\)，它通过过渡到商给出 \(u\)。
\end{enumerate}

\begin{enumerate}
\def\labelenumi{\alph{enumi})}
\item
  若 K 在 k 上形式平展，则前一习题的诸等价条件得到满足。（若 M 是 K-模，考察其底 k-向量空间为 K ⊕ M 的那个 k-代数，其中 M 是平方为零的理想。）
\item
  反之，若前一习题的诸等价条件得到满足，则 K 在 k 上形式平展。（由于域 K 在 k 上可分，习题 4 使我们能够证明 v 的存在性。）
\item
  设 \(\mathbf{B} = (b_{i})_{i\in\mathbf{I}}\) 是 K 中元素的族，使得 \((db_{i})_{i\in\mathbf{I}}\) 是 \(\Omega_{k}(\mathbf{K})\) 在 K 上的一个基。若 K 在 k 上可分，则 \(k(\mathbf{B})\) 是 k 的纯超越扩张（利用 A, V, 第 125 页定理 2，A, V, 第 97 页定理 1，以及 A, V, 第 165 页习题 6），且 K 在 \(k(\mathbf{B})\) 上形式平展。
\end{enumerate}

\begin{enumerate}
\def\labelenumi{\arabic{enumi})}
\setcounter{enumi}{8}
\tightlist
\item
  设 A 是等特征的局部环，\(m_{A}\) 是其极大理想，k 是 A 的一个子域。则剩余域 \(\kappa_{A}\) 是 k-代数。若 \(\kappa_{A}\) 在 k 上形式平展（习题 8），则称 k 是 A 的一个弱剩余域。
\end{enumerate}

\begin{enumerate}
\def\labelenumi{\alph{enumi})}
\item
  若 A 有一个弱表示域 \(k\)，则 A 的分离完备化 \(\hat{\mathbf{A}}\) 有唯一的表示域，它包含 \(k\) 在 \(\hat{\mathbf{A}}\) 中的像。
\item
  设 B 是包含于 A 中的局部环，其极大理想 \(m_{B} = m_{A} \cap B\)。若 \(\kappa_{A}\) 是 B 的剩余域 \(\kappa_{B}\) 的可分扩张，则 B 的每个弱剩余域都包含在 A 的一个弱剩余域中（利用习题 8 c)）。
\item
  设 B 是包含于 A 中的局部环，其极大理想 \(m_{B} = m_{A} \cap B\)。进一步假设 A 的特征 \(p\ne0\) 且 A 包含 \(B^{p}\)。则存在 A 的一个弱表示域，它包含 B 的一个弱表示域。
\end{enumerate}
% semantic-fragments: legacy-input-seam
\relax{}
